\documentclass[11pt,a4paper]{article}

% Japanese + Unicode support (XeLaTeX)
\usepackage{fontspec}
\usepackage{xeCJK}
\setmainfont{Times New Roman}
\setCJKmainfont{Hiragino Mincho ProN}
\setCJKmonofont{Hiragino Kaku Gothic ProN}

\usepackage{amsmath,amssymb}
\usepackage{graphicx}
\usepackage{booktabs}
\usepackage[margin=25mm]{geometry}
\usepackage{float}
\usepackage{hyperref}
\usepackage{subcaption}

\newcommand{\safeinclude}[2]{%
  \IfFileExists{#1}{%
    \includegraphics[width=#2]{#1}%
  }{%
    \fbox{%
      \parbox[c][35mm][c]{0.9\linewidth}{\centering Missing figure file}%
    }%
  }%
}

\title{N-Axis Attribution OoD (2026-02-16)}
\author{}
\date{2026-02-16}

\begin{document}
\maketitle

\section{目的}
本ドキュメントの目的は，\textbf{VLMにおけるAttribution OoDの軸設計フェーズ}を，
再現性・説明性・運用接続性の3観点で成立させることである．
ここでのAttribution OoDとは，単に「分布外かどうか」を判定するのではなく，
\textbf{どの種類の崩れが生じているか}を軸上で表現し，
後段の判定器や運用アクションに接続する設計である．

本稿の研究範囲は，Step Separationの原則に従い，
\textbf{最初$\rightarrow$軸作成}までに限定する．
したがって，後段のしきい値最適化や最終判定器の比較（AUROC最大化）を
同時に行わない．この切り分けにより，
「軸が良いから性能が上がったのか」
「判定器が良いから性能が上がったのか」
を識別できる状態を確保する．

本稿が答えるべき研究問いは以下の3つである．
\begin{enumerate}
  \item 軸数と正則化を，恣意的な重み関数なしで選べるか．
  \item 得られた軸はseed差分に対して安定か．
  \item 得られた軸を，運用上意味のある2軸（$z_u,z_c$）へ固定できるか．
\end{enumerate}

最終成果物は，固定軸の定義（数式），軸安定性評価，および可視化図である．
これらを後段判定器フェーズの固定入力設計として用いる．

\section{Attribution OoDの必要性(現場背景)}
従来のOoD判定は，入力をID/OoDの二値で扱う設計が中心であった．
しかし，VLMは事前学習規模が大きく，
「未知クラス」概念が従来より曖昧になるため，
二値判定のみでは運用に必要な情報が不足する．

現場で本当に必要なのは，OoD判定の\textbf{理由}である．
同じOoDでも，対処方法は大きく異なる．
例えば，
\begin{enumerate}
  \item 画像品質劣化（ブレ・ノイズ・低照度）なら再撮像で改善しやすい
  \item 意味的不一致（未知対象・文脈不整合）なら人手確認が必要
  \item 入力は明瞭だがモデル確信が低い場合は知識不足の可能性が高い
\end{enumerate}
となる．

この差を判定系へ渡すため，本研究では
\textbf{不確実性増大軸 $z_u$} と
\textbf{確信度低下軸 $z_c$} を分離して扱う．
これにより，運用アクションは次のように接続できる．

\begin{table}[H]
\centering
\begin{tabular}{lll}
\toprule
症状 & 軸空間での傾向 & 想定アクション \\
\midrule
入力品質劣化 & $z_u$高，$z_c$中〜低 & 再取得・再撮像・前処理改善 \\
意味的不一致 & $z_c$高，$z_u$中 & 人手介入・保留・ルール外処理 \\
複合異常 & $z_u$高，$z_c$高 & 強い警告と即時エスカレーション \\
通常入力 & $z_u$低，$z_c$低 & 自動処理継続 \\
\bottomrule
\end{tabular}
\caption{Attribution OoDの運用接続イメージ}
\label{tab:ops}
\end{table}

すなわち，本研究の目的は「精度指標を1点上げること」ではなく，
\textbf{誤りの原因を運用可能な座標系へ落とすこと}にある．

\section{背景と問題点}
本テーマの背景には，VLM時代のOoD検知における3つの構造的課題がある．

第1に，\textbf{軸数選択の恣意性}である．
多くの実装では，性能・安定性・解釈性を混ぜた合成スコアを使うが，
係数の選び方次第で最適解が変わる．
これは研究者自由度を増やし，再現性を下げる．

第2に，\textbf{小規模評価の不安定性}である．
サンプル数が小さいと，ロングテールの失敗パターンを十分に観測できない．
結果として，特定分割に依存した軸が選ばれやすくなる．

第3に，\textbf{軸定義の手動上書き}である．
SparsePCAで軸候補を得ても，最終的に手作業で軸式を作り直すと，
「データ駆動で軸を得た」という主張が弱くなる．

本研究では，これらを以下の方針で回避する．
\begin{enumerate}
  \item データ規模をID/OoD合計20,000まで拡大
  \item 分散閾値で定数特徴を除去し，数値不安定を先に抑止
  \item SparsePCAのCV再構成誤差に基づき1SEルールで選択
  \item 軸学習結果をそのまま固定し，後段と分離
\end{enumerate}

この設計により，「軸を先に固定してから判定器を評価する」
実験秩序を維持できる．

\section{解決方針(何を解くか)}
解く問題を次のように明確化する．

\textbf{入力}: VLMの推論ログから作る差分特徴 $\mathbf{x}_i$．

\textbf{出力}: 再現可能に固定された軸群 $\{z_j\}_{j=1}^{k}$ と，
運用用2軸 $(z_u,z_c)$ への割当．

\textbf{制約}: 軸選択は統計的手続き（CV+1SE）で行い，
最終判定器情報を使わない（情報リーク禁止）．

実装上の意思決定は以下である．
\begin{enumerate}
  \item 軸抽出法はSparsePCAを採用（線形・スパース・可読性重視）
  \item 軸数 $k$ と正則化 $\alpha$ はグリッド探索
  \item 評価指標は再構成誤差MSE（CV平均と標準誤差）
  \item 1SEルールで最小複雑モデルを採用
\end{enumerate}

この方針の利点は，
\textbf{「高性能そうに見える設定」ではなく「再現可能で説明可能な設定」}
を機械的に選べる点である．

\section{実行環境}
\begin{table}[H]
\centering
\begin{tabular}{ll}
\toprule
項目 & 内容 \\
\midrule
実行基盤 & Google Colab (Python 3.12) \\
GPU & NVIDIA A100-SXM4-40GB \\
PyTorch & 2.9.0+cu128 \\
主要ライブラリ & open-clip-torch 2.26.1, scikit-learn 1.5.2, pandas 2.2.x \\
VLM backend & OpenCLIP ViT-B-32 (laion2b\_s34b\_b79k) \\
ID dataset & TFDS food101 train[:10000] \\
OoD dataset 1 & TFDS cifar100 train[:5000] \\
OoD dataset 2 & TFDS imagenet\_r test[:5000] \\
総サンプル数 & 20,000 \\
出力成果物 & 軸定義，安定性評価，可視化図 \\
\bottomrule
\end{tabular}
\caption{軸作成実験の環境設定}
\label{tab:env}
\end{table}

補足として，初回実行時のTFDSダウンロード時間は長くなるが，
再実行時はローカルキャッシュを再利用するため高速化される．
また，CUDA系の警告（cuDNN/cuBLAS重複登録）は環境由来であり，
本実験の失敗原因ではなかった．

\section{処理フロー}
軸作成の手順を，入出力を含めて具体化する．

\begin{enumerate}
  \item \textbf{データ取得}\\
  TFDSからID/OoDを所定件数読み込む．

  \item \textbf{VLM推論}\\
  OpenCLIPで各サンプルのクラス分布を計算し，
  $\mathrm{conf}$，entropy，energy，ood\_scoreを記録する．

  \item \textbf{差分特徴構成}\\
  ID平均を基準に，drop/gain特徴を作る
  （確信度低下，不確実性増大など）．

  \item \textbf{前処理}\\
  Variance Thresholdで定数特徴を削除し，
  何を落としたかを実験ログとして記録する．

  \item \textbf{モデル選択}\\
  候補 $(k,\alpha)$ でCV再構成誤差を計算し，
  1SE選択用の比較表を得る．

  \item \textbf{1SE選択}\\
  最小MSEと標準誤差から閾値を作り，
  その範囲で最小複雑モデルを選ぶ．

  \item \textbf{最終軸固定}\\
  選択された $(k,\alpha)$ で全データ再学習し，
  軸荷重と軸スコアを固定する．

  \item \textbf{運用軸割当}\\
  不確実性寄与が最大の軸を $z_u$，確信度寄与が最大の軸を $z_c$ として固定する．
\end{enumerate}

このフローにより，後段判定器の入力は完全に固定される．

\section{数理定義}
\subsection*{サンプル指標}
サンプル $i$ の予測確率分布を $\mathbf{p}_i=(p_{i1},\ldots,p_{iC})$ とする．
最大確率を
\begin{equation}
\mathrm{conf}_i = \max_j p_{ij}
\end{equation}
と定義し，エントロピーを
\begin{equation}
H_i = -\sum_{j=1}^{C} p_{ij}\log(p_{ij}+\epsilon)
\end{equation}
とする（$\epsilon>0$は数値安定化）．

クラス数依存を除くため，正規化エントロピー
\begin{equation}
\tilde{H}_i = \frac{H_i}{\log(\max(2,C))}
\end{equation}
を用いる．
複合OoDスコアは
\begin{equation}
s_i^{\mathrm{ood}} = 0.5(1-\mathrm{conf}_i)+0.5\tilde{H}_i
\end{equation}
とする．

\subsection*{差分特徴}
ID集合を $\mathcal{D}_{\mathrm{ID}}$ とし，その平均を
$\bar{m}_{\mathrm{ID}}$ で表す．
例えば
\begin{equation}
d_{\mathrm{conf\_drop},i}=\bar{\mathrm{conf}}_{\mathrm{ID}}-\mathrm{conf}_i,
\end{equation}
\begin{equation}
d_{\mathrm{entropy\_gain},i}=\tilde{H}_i-\bar{\tilde{H}}_{\mathrm{ID}},
\end{equation}
\begin{equation}
d_{\mathrm{oodscore\_gain},i}=s_i^{\mathrm{ood}}-\bar{s}^{\mathrm{ood}}_{\mathrm{ID}}
\end{equation}
を定義する．
同様に $d_{\mathrm{msp\_drop}}, d_{\mathrm{energy\_gain}}, d_{\mathrm{abstain\_gain}}$ を作り，
\begin{equation}
\mathbf{x}_i=[d_{\mathrm{conf\_drop}}, d_{\mathrm{msp\_drop}}, d_{\mathrm{entropy\_gain}}, d_{\mathrm{energy\_gain}}, d_{\mathrm{oodscore\_gain}}, d_{\mathrm{abstain\_gain}}]^\top
\end{equation}
を軸学習の入力とする．

\subsection*{分散閾値前処理}
特徴 $x_j$ の標本分散を $\mathrm{Var}(x_j)$ とし，
\begin{equation}
\mathrm{Var}(x_j) < \tau_{\mathrm{var}}
\end{equation}
を満たす特徴を削除する（本実験では $\tau_{\mathrm{var}}=10^{-10}$）．
これにより，定数特徴由来の数値不安定を防ぐ．

\subsection*{SparsePCA}
標準化後行列を $X\in\mathbb{R}^{n\times d}$ とする．
SparsePCAは，成分行列 $W$ にL1型のスパース性を課し，
再構成誤差を最小化する方向を求める．
本実装ではscikit-learnの\texttt{SparsePCA}を用い，
候補 $k\in\{1,2,3,4\}$，
$\alpha\in\{0.5,1.0,2.0,4.0,8.0\}$ を探索した．

\subsection*{CV再構成誤差と1SEルール}
fold $f$ における検証誤差を
\begin{equation}
\mathrm{MSE}_f(k,\alpha)=\frac{1}{|\mathcal{V}_f|}\sum_{i\in\mathcal{V}_f}\lVert \mathbf{x}_i-\hat{\mathbf{x}}_i^{(k,\alpha)}\rVert_2^2
\end{equation}
とし，
\begin{equation}
\overline{\mathrm{MSE}}(k,\alpha)=\frac{1}{F}\sum_{f=1}^{F}\mathrm{MSE}_f(k,\alpha),
\quad
\mathrm{SE}(k,\alpha)=\frac{\mathrm{Std}(\mathrm{MSE}_f)}{\sqrt{F}}
\end{equation}
を計算する．

最小誤差を与える候補を $(k^{\ast},\alpha^{\ast})$ とし，
\begin{equation}
\tau_{\mathrm{1SE}}=\overline{\mathrm{MSE}}(k^{\ast},\alpha^{\ast})+\mathrm{SE}(k^{\ast},\alpha^{\ast})
\end{equation}
を定義する．
採用候補集合
\begin{equation}
\mathcal{C}=\{(k,\alpha)\mid \overline{\mathrm{MSE}}(k,\alpha)\le\tau_{\mathrm{1SE}}\}
\end{equation}
から，最小の $k$（同率なら大きい $\alpha$）を選ぶ．

\subsection*{最終固定軸}
本実験で選択された設定は
\begin{equation}
k=4,\quad \alpha=8.0
\end{equation}
である．
そのとき，運用軸として割り当てた式は
\begin{equation}
z_u = +0.9187\,d_{\mathrm{entropy\_gain}} +0.3949\,d_{\mathrm{oodscore\_gain}}
\end{equation}
\begin{equation}
z_c = +0.6546\,d_{\mathrm{conf\_drop}} +0.6546\,d_{\mathrm{msp\_drop}} +0.3783\,d_{\mathrm{oodscore\_gain}}
\end{equation}
である．
実装上は $\mathrm{msp}=\mathrm{conf}$ であるため，
\begin{equation}
z_c \simeq +1.3092\,d_{\mathrm{conf\_drop}} +0.3783\,d_{\mathrm{oodscore\_gain}}
\end{equation}
と等価整理できる．

\section{結果}
\subsection*{基礎統計}
ID基準平均は
\begin{equation}
\bar{\mathrm{conf}}_{\mathrm{ID}}=0.8259888770,
\quad
\bar{\tilde{H}}_{\mathrm{ID}}=0.1187297247,
\quad
\bar{s}^{\mathrm{ood}}_{\mathrm{ID}}=0.1463704238
\end{equation}
であった．

\subsection*{1SE選択結果}
\begin{table}[H]
\centering
\begin{tabular}{lcc}
\toprule
項目 & 値 & 解釈 \\
\midrule
$k_{\mathrm{selected}}$ & 4 & 4軸で十分な再構成 \\
$\alpha_{\mathrm{selected}}$ & 8.0 & 強めのスパース化 \\
$\mathrm{MSE}_{\min}$ & $8.3656\times 10^{-5}$ & 最小CV誤差 \\
$\mathrm{SE}_{\min}$ & $4.7763\times 10^{-7}$ & 誤差の推定ばらつき \\
$\tau_{\mathrm{1SE}}$ & $8.4134\times 10^{-5}$ & 許容閾値 \\
\bottomrule
\end{tabular}
\caption{1SEルールによる選択結果}
\label{tab:one_se}
\end{table}

\subsection*{seed安定性}
seed 42/43/44で再実行した結果，軸方向のコサイン類似度（絶対値）は
ほぼ1.0であった．

\begin{table}[H]
\centering
\begin{tabular}{lccccc}
\toprule
seed & $k$ & $\alpha$ & axis\_u & axis\_c & $|\cos|$（seed42基準）\\
\midrule
42 & 4 & 8.0 & 3 & 0 & $(z_u,z_c)=(1.000000,1.000000)$ \\
43 & 4 & 4.0 & 3 & 0 & $(z_u,z_c)=(0.999998,0.999999)$ \\
44 & 4 & 8.0 & 3 & 0 & $(z_u,z_c)=(1.000000,1.000000)$ \\
\midrule
平均 & - & - & - & - & $(0.99999948,0.99999983)$ \\
\bottomrule
\end{tabular}
\caption{3-seedでの軸安定性}
\label{tab:stability}
\end{table}

この結果は，
\textbf{軸の向きが乱数にほとんど依存しない}ことを示す．
後段で判定器のみを比較する条件として十分である．

\subsection*{図による検証結果}
図\\ref{fig:scatter}は，$z_u$-$z_c$ 空間でのID/OoD分布を直接示す．
図\\ref{fig:density}は，散布図の見かけ密度を補う等高線可視化であり，
両者の重なり領域と分離方向を明確にする．
図\\ref{fig:loadings}は，各軸の特徴寄与を定量的に示し，
式で定義した $z_u,z_c$ の意味づけを支える．
図\\ref{fig:cv}は，CV誤差曲線と1SE閾値を重ね，
$(k,\alpha)$ 選択が再現可能な手続きであることを可視化する．

\begin{figure}[H]
\centering
\safeinclude{figs/fig1_zu_zc_scatter.png}{0.8\linewidth}
\caption{$z_u$-$z_c$散布図（ID/OoD色分け）}
\label{fig:scatter}
\end{figure}

\begin{figure}[H]
\centering
\safeinclude{figs/fig2_zu_zc_density.png}{0.8\linewidth}
\caption{$z_u$-$z_c$密度等高線（ID/OoD）}
\label{fig:density}
\end{figure}

\begin{figure}[H]
\centering
\safeinclude{figs/fig3_axis_loadings.png}{0.95\linewidth}
\caption{SparsePCA軸荷重（axis loadings）}
\label{fig:loadings}
\end{figure}

\begin{figure}[H]
\centering
\safeinclude{figs/fig4_cv_1se.png}{0.85\linewidth}
\caption{CV再構成誤差と1SE選択}
\label{fig:cv}
\end{figure}

\section{可視化(グラフ)}
軸作成フェーズで必須の可視化は次の4種類である．

\begin{enumerate}
  \item \textbf{$z_u$-$z_c$散布図（ID/OoD色分け）}\\
  IDとOoDがどの方向で分離するかを視覚確認する．

  \item \textbf{2D密度等高線（KDE）}\\
  散布図だけでは見えにくい密度ピークと重なり領域を確認する．

  \item \textbf{軸loadings棒グラフ}\\
  各軸がどの特徴で構成されるかを定量確認する．

  \item \textbf{CV誤差曲線（1SE閾値付き）}\\
  選択した $(k,\alpha)$ の妥当性を説明する基礎図となる．
\end{enumerate}
本稿では，これら4種類を結果章に直接挿入し，
「数値表のみの主張」ではなく「視覚情報を伴う主張」に拡張した．

\section{考察}
本結果から得られる示唆は3点ある．

第1に，1SEルールの導入により，
軸数選択をヒューリスティック重み関数から切り離せた．
これは論文査読で問われる「なぜその軸数か」に対し，
手続き的かつ再現可能な回答を与える．

第2に，seed安定性が非常に高いことから，
データ規模20,000は軸固定フェーズとして妥当である．
この点は，先行の小規模検証で問題だった揺らぎを実質的に解消している．

第3に，$z_c$へ $d_{\mathrm{oodscore\_gain}}$ が寄与している点は重要である．
これは現実の失敗モードが「純粋な確信度低下」だけでなく，
不確実性要素を伴う混合モードである可能性を示す．
言い換えると，運用上は軸の完全直交を仮定しない方が安全である．

一方で制約も残る．
\begin{enumerate}
  \item OoD集合が自然画像寄り（cifar100, imagenet\_r）であり，
  ドメイン外シフトの種類は限定的である．
  \item プロンプト形式は固定であり，VLM特有のプロンプト感度は未評価である．
  \item 軸は線形結合であり，非線形失敗モードを直接は表現しない．
\end{enumerate}

したがって次工程では，軸固定を維持したまま，
判定器側で非線形境界や条件別評価を導入するのが妥当である．

\section{結論}
本稿では，研究計画の前半（最初$\rightarrow$軸作成）を完了し，
以下を達成した．

\begin{enumerate}
  \item ID/OoD合計20,000サンプルの軸作成基盤を構築
  \item 分散閾値 + SparsePCA + 1SEで恣意性の低い選択を実現
  \item $k=4,\alpha=8.0$ を選択し，運用軸 $z_u,z_c$ を確定
  \item 3-seedでほぼ完全一致の軸安定性を確認
\end{enumerate}

次の工程は，固定軸を入力にした判定器設計である．
このときもStep Separationを維持し，
判定器比較（AUROC/AUPR/FPR95，95\%信頼区間）を
軸の再学習なしで実施する．

以上より，本研究は「理由付きOoD検知」の前提条件である
軸定義の再現可能化を達成した．

\end{document}
